\documentclass[11pt]{article}
\usepackage[a4paper,margin=25mm]{geometry}
\usepackage{amsmath,amssymb,amsthm}
\usepackage[hidelinks]{hyperref}
\emergencystretch=2em
\newtheorem{theorem}{Theorem}
\title{Carleman Moments Do Not Recover an Unaugmented Path Law:\\A Disproof of Conjecture 00000004003 on Its Literal Process-Law Reading}
\author{GitHub: gaochengzhecpu}
\date{}
\begin{document}
\maketitle
\begin{abstract}
This counterexample concerns the original, unaugmented process law,
as stated in the bilingual source. It does not concern a law on
signatures or on tree-like equivalence classes. Two smooth based
one-dimensional closed paths have the same signature at every level,
but different laws as processes on $[0,1]$. Both laws have finite support
and satisfy an explicit divergent inverse-even-moment-root Carleman
condition. Lean verifies genuine recursive iterated integrals, all
signature levels, the moment condition, and inequality of the actual
path laws on the observation interval.
\end{abstract}

\paragraph{Reading and essential boundary.}
Both source languages say that the expected signature determines the
process distribution if and only if a Carleman condition holds. They
do not specify a tree-like quotient, time augmentation, or a
signature-injective class of paths. We refute the sufficiency direction
for the literal original path distribution. This boundary matters:
Hambly--Lyons identify signatures with tree-like equivalence classes
of finite-length paths \cite{HL}; Chevyrev--Lyons study probability
measures on signatures \cite{CL}. Our two paths are tree-like and have
the same signature law. Thus this example does not contradict either
of those determinacy frameworks. It does not address a corrected
source interpretation in which ``distribution'' means a signature
law or a law on tree-like equivalence classes, nor a time-augmented
signature. A Carleman hypothesis cannot restore information that the
unaugmented signature already discards deterministically.

\paragraph{The process family and its ordinary moment problem.}
For $a\in[1,2]$, set
\[
 x_a(t)=4at(1-t),\qquad 0\le t\le1.
\]
These are smooth paths with common starting point $0$ and common
endpoint $0$. Their midpoint values are $x_a(1/2)=a$, and
\[
 \|x_a\|_\infty=a,
\]
because $4t(1-t)=1-(2t-1)^2\in[0,1]$, with equality at $t=1/2$.
Consequently the family is injectively parametrized by amplitude,
recoverable from the original process. The ordinary amplitude moment
problem is a compactly supported scalar moment problem. We take
\[
 \mu=\delta_{x_1},\qquad \nu=\delta_{x_2}
\]
as Borel probability measures on $C([0,1],\mathbb R)$. They are
different: midpoint evaluation has law $\delta_1$ under $\mu$ and
$\delta_2$ under $\nu$.

\paragraph{Carleman convention and verification.}
The source does not specify the moment object in its Carleman
condition. We use the explicit radial path-norm version
\begin{equation}\label{eq:carleman}
 m_k(\eta)=\int\|p\|_\infty^k\,d\eta(p)<\infty,
 \qquad
 \sum_{n=1}^{\infty}m_{2n}(\eta)^{-1/(2n)}=\infty.
\end{equation}
For the family above this is exactly the ordinary positive-amplitude
Carleman condition. In particular,
\[
 m_{2n}(\mu)=1,\qquad m_{2n}(\nu)=2^{2n},
\]
so the respective summands in \eqref{eq:carleman} are $1$ and $1/2$.
All moments are finite and all even moments are strictly positive.

This choice does not exploit heavy tails or an indeterminate moment
problem. Any amplitude law supported in $[1,2]$ has inverse-even-root
terms at least $1/2$. For each of our two Dirac laws, every fixed
coordinate or finite collection of sampled coordinates is
deterministic and bounded, hence satisfies the usual scalar or
coordinatewise Carleman test, with the usual infinite inverse-root
convention for a zero moment. The total variation of $x_a$ is $2a$:
the path rises from $0$ to $a$ and falls back to $0$.
Thus a variation-norm version also has constant positive inverse-root
terms. We do not assert that there is one canonical infinite-dimensional
Carleman definition; these observations make the chosen convention
and the other common moment readings explicit.

\paragraph{Actual signatures at every level.}
For a smooth scalar path with $p(0)=0$, define its iterated integrals
recursively by
\[
 S_0(p;t)=1,\qquad
 S_{n+1}(p;t)=\int_0^t S_n(p;s)p'(s)\,ds.
\]
This is the one-dimensional tensor signature, with its unique scalar
coordinate at each tensor level. For our paths, induction gives
\begin{equation}\label{eq:signature}
 S_n(x_a;t)=\frac{x_a(t)^n}{n!}\qquad(n\ge0).
\end{equation}
The base case is $S_0=1$. If the formula holds at $n$, then
\[
 \frac{d}{dt}\frac{x_a(t)^{n+1}}{(n+1)!}
 =\frac{x_a(t)^n}{n!}x_a'(t).
\]
The fundamental theorem of calculus proves the next case, since
$x_a(0)^{n+1}=0$. All functions involved are polynomials and are
integrable on every compact interval. At $t=1$ the entire signature is
therefore
\[
 S(x_a)=(1,0,0,\ldots),
\]
independently of $a$.

\begin{theorem}
There are two distinct smooth based process laws on $[0,1]$ satisfying
\eqref{eq:carleman} whose expected signatures agree at every level.
Thus the source's claimed sufficiency fails for original unaugmented
process laws.
\end{theorem}
\begin{proof}
Use $\mu$ and $\nu$ above. The Carleman verification and inequality of
their midpoint laws were proved above. Integrating
\eqref{eq:signature} against either Dirac law gives the same expected
signature $(1,0,0,\ldots)$.
\end{proof}

Embedding the paths as $(x_a,0,\ldots,0)$ yields the same example in
every positive dimension. Any tensor word containing another
coordinate has zero iterated integral; the pure first-coordinate
words were computed above. Hence the obstruction is not confined
to dimension one. No separate assertion about an optimal truncation
level on a signature-injective class is made here.

\paragraph{Formal verification and scope.}
Lean defines $S_n$ by actual recursive Mathlib interval integrals
involving the derivative of a continuous real-valued path. It proves
\eqref{eq:signature} for every amplitude, level, and real upper endpoint
using derivative rules and the fundamental theorem of calculus.
It verifies smoothness, common starting point, midpoint recovery,
the exact supremum norm, two genuine Dirac probability measures,
integrability of all norm moments and signature coordinates, strict
positivity of even norm moments, and nonsummability of the actual
inverse-even-moment-root sequence. All expected-signature coordinates
are equal. The final theorem negates a universally quantified
Carleman-implies-determination proposition, including all these
admissibility assumptions.

For derivative bookkeeping, paths are stored as global smooth
extensions in $C(\mathbb R,\mathbb R)$. Crucially, the conclusion compares
their pushforward Borel laws on $C([0,1],\mathbb R)$, and formal
inequality is detected at the midpoint inside that interval. No
disagreement outside the observation interval is used.
The general-amplitude-law remark, variation-norm observation,
higher-dimensional embedding, and tree-like interpretation are proved
or explained here on paper rather than formalized. The formal result
is the explicit scalar counterexample with the stated radial Carleman
condition, not a refutation of a signature-law uniqueness theorem.

\paragraph{Source and reproduction.}
The exact bilingual source is \texttt{SOURCE.md}; its frozen commit
and SHA-256 are recorded in
\texttt{verification/SOURCE\_PROVENANCE.md}.
Use Lean 4.19.0 and the pinned Mathlib manifest. From \texttt{lean/}, run
\begin{center}\texttt{lake env lean -DwarningAsError=true Main.lean}\end{center}
The strict log records the real compiler result. The final axiom audit
uses only \texttt{propext}, \texttt{Classical.choice}, and
\texttt{Quot.sound}. Development was by a delegated AI agent with
self-review; the coordinating agent reviews it separately. No external
independent review is claimed.

\begin{thebibliography}{2}
\bibitem{HL} B. Hambly and T. Lyons,
\emph{Uniqueness for the signature of a path of bounded variation
and the reduced path group},
\href{https://arxiv.org/abs/math/0507536}{arXiv:math/0507536}.
\bibitem{CL} I. Chevyrev and T. Lyons,
\emph{Characteristic functions of measures on geometric rough paths},
\href{https://arxiv.org/abs/1307.3580}{arXiv:1307.3580}.
\end{thebibliography}
\end{document}
